% !TEX encoding = UTF-8 Unicode
% !TeX program = xelatex
% !TEX TS-program = XeLaTeX

\documentclass[10pt, a5paper]{article}
\usepackage{mathtools}
\usepackage{unicode-math}
\usepackage{comment}
\usepackage{tikz}
\usepackage[pass]{geometry}
%\geometry{legalpaper, landscape, margin=1in}
\usepackage{fullpage}
\usepackage{setspace}
\usepackage[localise]{xepersian}
\setstretch{1.3}
\settextfont[Scale=1, BoldFont=XB NiloofarBd.ttf]{XB Niloofar.ttf}
\setlatintextfont[Scale=1]{XB Niloofar}%{Junicode} %Junicode
\setdigitfont[BoldFont=XB NiloofarBd.ttf]{XB Niloofar.ttf}%{XB Zar}

\setlength{\parindent}{0ex}

\begin{document}
	
	\begin{tikzpicture}[remember picture,overlay]
	\node[anchor=north,yshift=-1.5cm] at (current page.north) {\rl{به نام خدا}};
	\end{tikzpicture}%

\vspace{-1cm}
{\centering
	\textbf{کوئیز اول درس ساختمان داده‌ها و الگوریتم‌ها}\\
مدرس: مهدی صفرنژاد       - ۱۵ مهرماه ۹۸\\
\raggedleft\small
نام و نام خانوادگی:\\
شماره دانشجویی:\\
}
\vspace{-0.5cm}
\پرش‌متوسط\پرخط‌ا\\
تمام سؤال‌ها جواب کوتاه هستند و توضیح یک یا دو خطی می‌خواهند.

\پاراگراف{سؤال ۱. }
رشد تابع زیر را به صورت مجانبی تحلیل کنید.
\vspace{-0.9cm}
\begin{latin}
\begin{align*}
T(n) = 4T(\sqrt{n}) + \log^2(n)
\end{align*}
\end{latin}


\vspace{3.7cm}

\پاراگراف{سؤال ۲.}
توابع داده شده در زیر را به ترتیبی مرتب کنید که تابع

$f_i \in\Omega(f_{i+1})$
باشد.
(
نیازی به نوشتن راه‌حل نیست!
)
\begin{gather*}
    (\log{n})^{\log{n}}, 1+2+...+n, (\sqrt{2})^n, (n)^{\log{n}}
\end{gather*}



\vspace{3cm}

\پاراگراف{سؤال ۳.}
اگر بخواهیم لیست
$
n
$
نفره‌ی دانشجویان کلاس را بر حسب شماره دانش‌جویی‌شان مرتب کنیم
و بدانیم که هر کس حداکثر ۵ خانه از جای اصلی‌اش در لیست فاصله دارد بهتر است از کدام الگوریتم مرتب‌سازی بهره بگیریم؟
مرتبه‌ی زمانی الگوریتم انتخابی در این صورت چه خواهد بود؟
\vspace{3cm}

% \پاراگراف{سؤال ۴.}
% C


\begin{flushleft}
شاد باشین :)
\end{flushleft}

\end{document}